%% ****** Start of file apstemplate.tex ****** %
%%
%%
%%   This file is part of the APS files in the REVTeX 4.2 distribution.
%%   Version 4.2a of REVTeX, January, 2015
%%
%%
%%   Copyright (c) 2015 The American Physical Society.
%%
%%   See the REVTeX 4 README file for restrictions and more information.
%%
%
% This is a template for producing manuscripts for use with REVTEX 4.2
% Copy this file to another name and then work on that file.
% That way, you always have this original template file to use.
%
% Group addresses by affiliation; use superscriptaddress for long
% author lists, or if there are many overlapping affiliations.
% For Phys. Rev. appearance, change preprint to twocolumn.
% Choose pra, prb, prc, prd, pre, prl, prstab, prstper, or rmp for journal
%  Add 'draft' option to mark overfull boxes with black boxes
%  Add 'showkeys' option to make keywords appear
%\documentclass[aps,prl,preprint,groupedaddress]{revtex4-2}
\documentclass[aps,prl,twocolumn,groupedaddress]{revtex4-2}
%\documentclass[aps,prl,preprint,superscriptaddress]{revtex4-2}
%\documentclass[aps,prl,reprint,groupedaddress]{revtex4-2}

% You should use BibTeX and apsrev.bst for references
% Choosing a journal automatically selects the correct APS
% BibTeX style file (bst file), so only uncomment the line
% below if necessary.
%\bibliographystyle{apsrev4-2}

\usepackage{graphicx}
\usepackage{epstopdf}
\usepackage{caption}
%\usepackage{amsmath}% http://ctan.org/pkg/amsmath
%\usepackage{amsthm}
%\usepackage{amsfonts}
%\usepackage{subfigure}
%\usepackage{hhline}
%\usepackage[miktex]{gnuplottex}
%\usepackage{xcolor}
\usepackage{amssymb}
\usepackage{amsmath}
\usepackage{color}
\usepackage{hyperref}
%\usepackage[percent]{overpic}
\usepackage{tikz}
\usepackage{mathrsfs}
\usepackage{wasysym}
\usepackage{tikz-cd}
%\usepackage{linearb}
%\usepackage{wsuipa}
\usepackage{calc}
\usepackage{tipa}
\usepackage{booktabs}

%\usepackage{stix} %\fisheye
\usepackage{stackengine,scalerel}

% so sections, subsections, etc. become numerated.
\setcounter{secnumdepth}{3}

\newcommand{\FF}{\mathbb{F}}
\newcommand{\NN}{\mathbb{N}}
\newcommand{\ZZ}{\mathbb{Z}}
\newcommand{\QQ}{\mathbb{Q}}
\newcommand{\RR}{\mathbb{R}}
\newcommand{\CC}{\mathbb{C}}
\newcommand{\II}{\mathbb{I}}
\newcommand{\GG}{\mathbb{G}}

\DeclareMathOperator*{\argmax}{arg\,max}
\DeclareMathOperator*{\argmin}{arg\,min}
\newcommand{\avrg}[1]{\left\langle #1 \right\rangle}
\newcommand{\nelta}{\bar{\delta}}
\newcommand{\bra}[1]{\left\langle #1\right|}
\newcommand{\ket}[1]{\left| #1 \right\rangle}
\newcommand{\sbra}[1]{\langle #1|}
\newcommand{\sket}[1]{| #1 \rangle}
\newcommand{\bek}[3]{\left\langle #1 \right| #2 \left| #3 \right\rangle}
\newcommand{\sbek}[3]{\langle #1 | #2 | #3 \rangle}
\newcommand{\braket}[2]{\left\langle #1 \middle| #2 \right\rangle}
\newcommand{\ketbra}[2]{\left| #1 \middle\rangle \middle\langle #2  \right|}
\newcommand{\sbraket}[2]{\langle #1 | #2 \rangle}
\newcommand{\sketbra}[2]{| #1 \rangle  \langle #2 |}
\newcommand{\norm}[1]{\left\lVert#1\right\rVert}
\newcommand{\snorm}[1]{\lVert#1\rVert}
\newcommand{\bvec}[1]{\boldsymbol{\mathsf{#1}}}
\newcommand{\bcov}[1]{\boldsymbol{#1}}
\newcommand{\bdua}[1]{\boldsymbol{\check{#1}}}
%\newcommand{\bdov}[1]{\boldsymbol{\breve{#1}}}
\newcommand{\bdov}[1]{\breve{#1}}
%\newcommand{\bten}[1]{\boldsymbol{\mathfrak{#1}}}
\newcommand{\bten}[1]{\boldsymbol{\mathfrak{#1}}}
\newcommand{\forany}{\tilde{\forall}}
\newcommand{\qed}{$\overset{\circ}{.}\;$}
\newcommand{\bigo}{\mathcal{O}}
\newcommand{\spn}{\mathrm{spn}\,}
\newcommand{\krn}{\mathrm{krn}\,}
\newcommand{\biy}{\leftrightarrow}
\newcommand{\ugh}{\mathbin{\textlinb{\BPwheel}}}

\newcommand\bigeye{\ensurestackMath{\stackinset{c}{}{c}{-.3pt}%
  {\bullet}{\scriptstyle\bigcirc}}}
\newcommand\eye{\scalerel*{\bigeye}{x}}
%\newcommand*{\fisheye}{%
%    \mathbin{%
%        \ooalign{$\circledcirc$\cr\hidewidth$\bullet$\hidewidth}%
%    }%
%}

% double angle-bracket
% https://tex.stackexchange.com/questions/79657/how-to-get-double-angle-bracket-without-using-mnsymbol-package
\makeatletter
\newsavebox{\@brx}
\newcommand{\llangle}[1][]{\savebox{\@brx}{\(\m@th{#1\langle}\)}%
  \mathopen{\copy\@brx\mkern2mu\kern-0.9\wd\@brx\usebox{\@brx}}}
\newcommand{\rrangle}[1][]{\savebox{\@brx}{\(\m@th{#1\rangle}\)}%
  \mathclose{\copy\@brx\mkern2mu\kern-0.9\wd\@brx\usebox{\@brx}}}
\makeatother

% https://tex.stackexchange.com/questions/297362/bar-through-letter-in-mathmode
%\usepackage{calc}
\newsavebox\CBox
\newcommand\hcancel[2][0.5pt]{%
  \ifmmode\sbox\CBox{$#2$}\else\sbox\CBox{#2}\fi%
  \makebox[0pt][l]{\usebox\CBox}%  
  \rule[0.77\ht\CBox-#1/2]{\wd\CBox}{#1}}
\newcommand{\cb}{\hcancel{b}}
\newcommand{\cd}{\hcancel{d}}
\newcommand{\ct}{\hcancel{\partial}}
%\newcommand{\supp}{\mathsf{supp}}
\newcommand{\supp}{\mathrm{supp}}
\newcommand{\csupp}{\overline{\mathrm{supp}}}
\newcommand{\supr}[2]{\underset{#1}{\text{supr}}}

\begin{document}

% Use the \preprint command to place your local institutional report
% number in the upper righthand corner of the title page in preprint mode.
% Multiple \preprint commands are allowed.
% Use the 'preprintnumbers' class option to override journal defaults
% to display numbers if necessary
%\preprint{}

%Title of paper
\title{
M\'etodos Matem\'aticos de la F\'isica II\\
Cheat Sheet: 
Differential Geometry
}

% repeat the \author .. \affiliation  etc. as needed
% \email, \thanks, \homepage, \altaffiliation all apply to the current
% author. Explanatory text should go in the []'s, actual e-mail
% address or url should go in the {}'s for \email and \homepage.
% Please use the appropriate macro foreach each type of information

% \affiliation command applies to all authors since the last
% \affiliation command. The \affiliation command should follow the
% other information
% \affiliation can be followed by \email, \homepage, \thanks as well.
\author{Juan I. Perotti}
\email[]{juan.perotti@unc.edu.ar}
%\homepage[]{Your web page}
%\thanks{}
%\altaffiliation{}
%\affiliation{}
\affiliation{Instituto de F\'isica Enrique Gaviola (IFEG-CONICET), Ciudad Universitaria, 5000 C\'ordoba, Argentina}
\affiliation{Facultad de Matem\'atica, Astronom\'ia, F\'isica y Computaci\'on, Universidad Nacional de C\'ordoba, Ciudad Universitaria, 5000 C\'ordoba, Argentina}

% \author{Colaborador}
% \affiliation{Instituto de F\'isica Enrique Gaviola (IFEG-CONICET), Ciudad Universitaria, 5000 C\'ordoba, Argentina}
% \affiliation{Facultad de Matem\'atica, Astronom\'ia, F\'isica y Computaci\'on, Universidad Nacional de C\'ordoba, Ciudad Universitaria, 5000 C\'ordoba, Argentina}
% \email[E-mail: ]{xxx.yyy@unc.edu.ar }

%Collaboration name if desired (requires use of superscriptaddress
%option in \documentclass). \noaffiliation is required (may also be
%used with the \author command).
%\collaboration can be followed by \email, \homepage, \thanks as well.
%\collaboration{Claudio J. Tessone}
%\noaffiliation

\date{\today}

%\begin{abstract}
%Bla bla bla... resumen...
%\end{abstract}

% insert suggested keywords - APS authors don't need to do this
%\keywords{}

%\maketitle must follow title, authors, abstract, and keywords
\maketitle
\onecolumngrid

\section{Manifolds}

An $n$-dimensional manifold $M$ is a set equipped with a set of invertible maps $\varphi_i:D_i\leftrightarrow \RR^n$ between open subsets $D_i$ of $M$ and $\RR^n$.
Crucially, this set of maps is defined to be such that $\cup_i D_i = M$, i.e. the union of the domains of these maps cover $M$.
Moreover, these maps are required to satisfy a very important property.
Namely, if $p \in D_i\cap D_j$, i.e. it is a point in the intersection of two of such domains, and $x = \varphi_i(p)\in \RR^n$, then the map
$$
\varphi_j \circ \varphi_i^{-1} : \RR^n \leftrightarrow D_i\cap D_j \leftrightarrow \RR^n
$$
is infinitely differentiable at $x$.
This is how the notions of {\em continuity} and {\em smoothness} is attributed to $M$.
A set of such maps is called an {\em atlas} of $M$.
Any of such maps $\varphi_i$ is called a {\em chart}.

Remark: often a chart $\varphi_i$ is defined by mapping $D_i$ to an open subset $U$ of $\RR^n$ as opposed to the full euclidean space $\RR^n$, which looks more general.
However, since we can allways find an invertible smooth map between any open subset $U$ of $\RR^n$ and $\RR^n$ itself, our definitions are as general as the alternative.

\subsection{Coordinate systems}

Any chart $\varphi:D\leftrightarrow M$ defines a {\em coordinate system}, assigning the coordinates
$$
\RR^n \ni (x^1,...,x^n) = x = \varphi(p)
$$
to each point $p\in D$.
We can always add new maps to an atlas, given that the prescribed properties are satisfied.
This is how we can always define new coordinate systems over some open domain $D$ of $M$.

\section{Tangent spaces}

\subsection{Derivations}

Let $C^{\infty}(M)$ be the set of {\em smooth} functions (i.e. infinitely differentiable) from $M$ to $\RR$.

A {\em derivation} over a given and fixed point $p$ in $M$ is a map $v:C^{\infty}\to \RR$ such that it satisfies:
\begin{itemize}
\item linearity
$$
v(a\,f+b\,g) = a\,v(f)+b\,v(g),
$$
\item and Leibniz rule
$$
v(fg)
=
f(p)\,v(g)
+
g(p)\,v(f)
$$
\end{itemize}
for any two arbitrary smooth functions $f,g\in C^{\infty}(M)$ and any two constants $a,b\in \RR$.
Here, $fg$ denotes the function such that
$$
(fg)(p)
=
f(p)\,g(p)
$$
for each $p\in M$.

\subsection{Tangent space}

The set of derivations over a given point $p\in M$ is a vector space called the {\em tangent space} of $M$ at $p$, and is denoted by $T_pM$.
The elements or vectors in $T_pM$ are called tangent vectors.

\subsection{Coordinate basis}

Given a coordinate system $p \to \varphi(p) = (x^1,...,x^n)$, a partial derivative operator $\partial_i$ defined by
$$
(\partial_i f)
:=
\frac{\partial (f\circ \varphi^{-1})}{\partial x^i}(x=\varphi(p))
$$
for any function $f\in C^{\infty}(M)$ and a given and fixed $p=\varphi^{-1}(x) \in M$, is a tangent vector in $T_pM$.
Namely, $\partial_i$ is a derivation at $p$ because, it is clearly linear and it satisfies Leibniz rule
\begin{eqnarray}
(\partial_i (fg))
&=&
\frac{\partial (fg\circ \varphi^{-1})}{\partial x^i}(x=\varphi(p))
\notag
\\
&=&
\frac{\partial (f\circ \varphi^{-1}) (g\circ \varphi^{-1})}{\partial x^i}(x=\varphi(p))
\notag
\\
&=&
(g\circ \varphi^{-1} \circ \varphi)(p)
\,
\frac{\partial (f\circ \varphi^{-1})}{\partial x^i}(x=\varphi(p))
+
(f\circ \varphi^{-1} \circ \varphi)(p)
\,
\frac{\partial (g\circ \varphi^{-1})}{\partial x^i}(x=\varphi(p))
\notag
\\
&=&
g(p)
\,
(\partial_i f)
+
f(p)
\,
(\partial_i g)
\notag
\end{eqnarray}
for any smooth functions $f,g\in C^{\infty}(M)$.

Crucially, without proof, we claim that the set of partial derivatives $\{\partial_1,...,\partial_n\}$ is a basis of $T_pM$ called the coordinate basis associated to the coordinate system defined by a chart $\varphi:D\ni p\to \varphi(p)=(x^1,...,x^n)\in \RR^n$.
Hence, any tangent vector $v\in T_pM$ can be written as the linear combination
$$
v = v^i\partial_i,
$$
where the $v^i\in \RR$ are the unique scalars or coordinates representing $v$.

\section{Differential 1-forms}

Let $T_p^*M$ denote the dual space of $T_pM$.
If $\{\partial_1,...,\partial_n\}$ is a coordinate basis of $T_pM$, let $\{dx^1,...,dx^n\}$ denote the corresponding dual basis, so
$$
dx^i(\partial_j)
=
\delta^i{}_j.
$$
The covectors or dual basis vectors $dx^i$ are called {\em differential 1-forms} or simply 1-forms.

\appendix

\section{Coordinate basis}

Here we show that $\{\partial_1,...,\partial_n\}$ is a basis of $T_pM$.

\noindent {\bf 1. Tangent vectors as derivations}

Let $M$ be an $n$-dimensional smooth manifold and let $p\in M$.
The tangent space $T_pM$ is defined as the set of linear maps
$$
X:C^\infty(M)\longrightarrow\mathbb{R}
$$
satisfying the Leibniz rule
$$
X(fg)=X(f)g(p)+f(p)X(g).
$$
Such maps are called {\em derivations at $p$}.

\noindent {\bf 2. Coordinate functions}

Choose a coordinate chart
$$
(x^1,\ldots,x^n)
$$
defined near $p$.
Each coordinate function
$$
x^i:M\longrightarrow\mathbb{R}
$$
is smooth.

\noindent {\bf 3. Coordinate derivations}

For every $i$, define
$$
\boxed{
\partial_i(f)
:=
\left.
\frac{\partial(f\circ x^{-1})}{\partial x^i}
\right|_{x(p)}.
}
$$
That is,
\begin{enumerate}
\item express $f$ in coordinates,
\item take the ordinary partial derivative,
\item evaluate at the coordinates of $p$.
\end{enumerate}
These are derivations because ordinary partial derivatives satisfy linearity and the product rule.

\noindent {\bf 4. Linear independence}

Suppose
$$
a^i\partial_i=0.
$$
We apply this derivation to the coordinate functions.
Since
$$
\partial_i(x^j)=\delta_i^j,
$$
we obtain
$$
0
=
\left(a^i\partial_i\right)(x^j)
=
a^i\delta_i^j
=
a^j.
$$
Thus every coefficient vanishes.
Hence
$$
\boxed{
\{\partial_1,\ldots,\partial_n\}
\text{ is linearly independent.}
}
$$

\noindent {\bf 5. Spanning}

Now let
$$
X\in T_pM
$$
be an arbitrary derivation.
We want to show
$$
X
=
X^i\partial_i.
$$
The key is to determine the coefficients.
Define
$$
X^i:=X(x^i).
$$
These are just real numbers.

Now consider an arbitrary smooth function $f$ over $M$.

\begin{itemize}
\item Taylor expansion.

Near $p$,
$$
f(x)
=
f(p)
+
(x^i-x^i(p))
\frac{\partial f}{\partial x^i}(p)
+
R(x),
$$
where
$$
R(x)=O(|x-x(p)|^2).
$$

\item Apply the derivation.

Since derivations kill constants,
$$
X(f(p))=0.
$$
Therefore
$$
X(f)
=
X
\left(
(x^i-x^i(p))
\frac{\partial f}{\partial x^i}(p)
+
R
\right).
$$
Because
$$
\frac{\partial f}{\partial x^i}(p)
$$
is just a constant,
$$
X(f)
=
\frac{\partial f}{\partial x^i}(p)
X(x^i)
+
X(R).
$$
Thus everything boils down to showing
$$
X(R)=0.
$$

\item Why $X(R)=0$.

Since $R$ vanishes to second order,
it can be written locally as
$$
R
=
\sum_{i,j}
(x^i-x^i(p))
(x^j-x^j(p))
r_{ij}(x),
$$
for smooth functions $r_{ij}$.
Applying the Leibniz rule,
every term contains a factor
$$
x^k-x^k(p),
$$
which vanishes at $p$.
For example,
$$
\begin{aligned}
X\left((x^i-a^i)(x^j-a^j)\right)
&=
X(x^i-a^i)(x^j-a^j)(p)\\
&\quad+
(x^i-a^i)(p)X(x^j-a^j)\\
&=0.
\end{aligned}
$$
Hence
$$
X(R)=0.
$$
\end{itemize}
Therefore, 
$$
X(f)
=
X(x^i)
\frac{\partial f}{\partial x^i}(p).
$$
Since
$$
\partial_i(f)
=
\frac{\partial f}{\partial x^i}(p),
$$
we obtain
$$
\boxed{
X
=
X(x^i)\partial_i.
}
$$
Thus every derivation is a linear combination of the coordinate derivations.


\noindent {\bf 6. Conclusion}

Since the $\partial_i$ are linearly independent and span $T_pM$,
$$
\boxed{
\left\{
\frac{\partial}{\partial x^1},
\ldots,
\frac{\partial}{\partial x^n}
\right\}
\text{ is a basis of }T_pM.
}
$$

\noindent {\bf Geometric interpretation}

This proof also explains {\em why the coordinate functions are so special}. A derivation $X$ is completely determined by its action on the $n$ coordinate functions:
$$
X
\longleftrightarrow
\bigl(X(x^1),\ldots,X(x^n)\bigr)\in\mathbb{R}^n.
$$
Indeed, the map
$$
X\mapsto \bigl(X(x^1),\ldots,X(x^n)\bigr)
$$
is an isomorphism $T_pM\simeq\mathbb{R}^n$. In coordinates, a tangent vector is therefore identified with its components $X^i=X(x^i)$, and the basis vectors $\partial_i$ are precisely those derivations that satisfy
$$
\partial_i(x^j)=\delta_i^{,j}.
$$
This is the derivation-based analogue of the standard basis $(1,0,\ldots,0),\ldots,(0,\ldots,0,1)$ of $\mathbb{R}^n$. It also makes clear why tangent vectors are often thought of as ``directional derivatives'': once a coordinate system is chosen, every derivation acts exactly like a first-order differential operator
$$
X = X^i\frac{\partial}{\partial x^i}.
$$

\end{document}
%
% ****** End of file apstemplate.tex ******


